This chapter reviews the research most relevant to MPVRP-CC. It brings
together two lines of work: routing models that account for repeated trips,
multiple products, and preparation activities; and solution methods that
combine routing decisions with constraint programming. The purpose is to
identify the foundations of this study and explain how its questions relate
to existing contributions.

The review is thematic rather than exhaustive. It draws on selected routing
surveys, application studies, and papers on CP representations. The comparison
focuses on the decisions being optimized, the operational meaning of cleaning
or setup, and the role of the solution method. These distinctions matter
because similar terms can describe different mathematical problems.

\section{Relevant routing variants}\label{sec:routing-variants}

\subsection{Multiple trips}

In a multi-trip routing problem, a vehicle can carry out more than one trip
within the planning horizon. Cattaruzza, Absi, and
Feillet~\cite{cattaruzza2016multipletrips} survey this family, its formulations,
and its exact and heuristic methods. They also examine related routing
settings in which repeated vehicle use is allowed. Their review shows that
multiple trips are a structural feature shared by several application areas,
rather than a single narrowly defined problem.

For the present study, the important consequence is that trips assigned to
the same vehicle cannot always be evaluated independently. Once the previous
load affects the preparation for the next one, the order of trips becomes
part of the cost calculation. Multi-trip routing therefore provides the
setting for the changeover question, but allowing several trips does not
by itself specify how product transitions should be charged.

\subsection{Multiple products and operational constraints}

Multi-product routing introduces distinctions between the goods being
transported. These distinctions can affect loading, compatibility, and the
allocation of vehicle capacity. Multi-compartment routing is one way to
represent the simultaneous transport of products that must remain separate;
the olive-oil application of Lahyani et
al.~\cite{lahyani2015oilmulticompartment} illustrates this connection.
However, carrying several products in separate compartments and carrying
one product after another are different decisions. MPVRP-CC emphasizes the
latter interaction between successive loads.

Other features concern the resources available to serve demand. Vehicles
may have different capacities, several depots may supply products, and a
customer's demand may be divided between deliveries. The broader routing
literature discusses such extensions and the ways they change the
structure of a plan~\cite{toth2014vehicle,wen2010rich}. For MPVRP-CC, their
importance lies in their interaction: the available stock can influence
where a vehicle loads, capacity can influence how deliveries are grouped,
and those choices can affect which product transitions occur.

A useful distinction is between a vehicle's home location and a loading
depot. Assigning a vehicle to a home garage does not necessarily determine
where it obtains every load. Likewise, the presence of several products does
not imply that every station requires all of them. These details must be
specified when comparing models; otherwise, two descriptions that both use
``multi-product'' or ``multi-depot'' may conceal different decision spaces.
The formal conventions adopted in this thesis are given in
Section~\ref{sec:problem-definition}.

\subsection{Rich vehicle routing problems}

Lahyani, Khemakhem, and Semet~\cite{lahyani2015taxonomy} organize rich vehicle
routing research through a taxonomy and a discussion of its definition.
Wen's thesis~\cite{wen2010rich} also addresses rich routing applications.
These works provide a framework for discussing models that bring together
several operational features, rather than studying each extension in
isolation.

MPVRP-CC belongs within this broader perspective because its routing,
loading, capacity, stock, and product-sequencing decisions interact. This
classification describes its scope; it does not establish novelty. The
relevant comparison is whether existing studies already represent the same
preparation mechanism and the same set of operational decisions.

\section{Cleaning and changeover costs}\label{sec:cleaning-changeover}

\subsection{Sequence-dependent preparation}

The idea of making preparation depend on the preceding activity is well
established outside routing. Haase~\cite{haase1994setup}, for example, studies
capacitated lot-sizing with sequence-dependent setup costs. The connection
to this thesis is the dependence on order: the cost of starting an activity
can depend on the activity completed immediately before it.

Transferring this idea to routing requires specifying what is prepared and
when. A production setup changes the state of a production resource; a
vehicle changeover concerns the preparation for another load. A monetary
charge must also be distinguished from a preparation duration or a
prohibited transition. These mechanisms may coexist, but they affect the
model in different ways. A cost changes the objective, a duration can affect
scheduling feasibility, and an incompatibility excludes an operation unless
an appropriate preparation action is performed.

\subsection{Vehicle and compartment cleaning}

Lahyani et al.~\cite{lahyani2015oilmulticompartment} study multi-period
olive-oil collection with compartmented vehicles. A compartment carrying
the lowest grade requires cleaning before a superior grade is loaded, and
the formulation includes a cleaning charge. This is a close precedent:
product history affects preparation on the transport resource itself.
MPVRP-CC instead concerns deliveries and product-dependent costs between
successive loads. The comparison therefore concerns the operating structure
and transition rule, not the introduction of cleaning into routing for the
first time.

Chokanat, Pitakaso, and Sethanan~\cite{agriengineering1010006} study raw-milk
transportation with an objective that includes transportation and truck and
tank cleaning costs. They propose a modified differential evolution method.
The study is relevant because cleaning is part of the economic evaluation
of a transport plan. However, the presence of a cleaning charge alone does
not establish that the charge depends on an ordered pair of different
products. It is therefore treated here as a cleaning-cost precedent, rather
than as evidence of an identical changeover mechanism.

\subsection{Cleaning in maritime routing}

Tamburini, Lange, and Pisinger~\cite{TAMBURINI2025104019} study tramp ship
routing and scheduling in a setting with contract and spot cargoes. Their
model combines several operational decisions, including speed, chartering,
bunkering, and cleaning, and is solved using column generation. It extends
the comparison beyond road transport and illustrates the integration of
preparation activities into a rich routing and scheduling problem.

This connection is important, but similarity at the level of an application
description does not establish equivalence between formulations. A cleaning
state or a particular cleaning rule should not automatically be interpreted
as an arbitrary product-to-product cost matrix. The present comparison
therefore identifies this work as a related integrated model without
asserting that it is either identical to, or unable to represent, MPVRP-CC.

\subsection{Production setups coupled with routing}

Miranda et al.~\cite{MIRANDA2018211} investigate a rich production-routing
problem motivated by a furniture manufacturer. It includes multiple
products, sequence-dependent setup times, a heterogeneous fleet, and delivery
requirements. Their decomposition heuristic iterates between a subproblem
for production and customer assignment and a routing subproblem, with
information from routing used to revise the first subproblem.

Here, the sequence-dependent preparation takes place on the production
line. Its effect on distribution is transmitted through the coordination of
production and delivery. In MPVRP-CC, the preparation concerns the vehicle
between loads. This difference identifies the resource whose history must
be represented. The study nevertheless offers a useful methodological
comparison because it addresses interacting decisions through decomposition.

\subsection{Synthesis}

Table~\ref{tab:cleaning-literature} summarizes the mechanisms relevant to this
review. It is a comparison of selected features, not a claim that the papers
have otherwise equivalent assumptions. In particular, a study can be closely
related through its preparation mechanism while differing substantially in
its demand, loading, or scheduling structure.

\begin{table}[htbp]
  \centering
  \small
  \caption{Preparation mechanisms in selected related studies}
  \label{tab:cleaning-literature}
  \begin{tabularx}{\textwidth}{@{}p{0.22\textwidth}XX@{}}
    \toprule
    Study & Relevant mechanism & Main distinction for this thesis \\
    \midrule
    Lahyani et al.~\cite{lahyani2015oilmulticompartment}
      & Compartment cleaning linked to previously carried oil grade
      & Multi-period collection with compartments \\[0.5em]
    Chokanat et al.~\cite{agriengineering1010006}
      & Truck and tank cleaning charges
      & Pairwise product dependence not established in this comparison \\[0.5em]
    Tamburini et al.~\cite{TAMBURINI2025104019}
      & Cleaning integrated with ship routing and scheduling
      & Maritime operating structure; model equivalence not asserted \\[0.5em]
    Miranda et al.~\cite{MIRANDA2018211}
      & Sequence-dependent production setup times
      & Preparation occurs on the production line \\
    \bottomrule
  \end{tabularx}
\end{table}

The common lesson is that preparation can be part of an optimization model
rather than a fixed expense added after planning. The remaining task is to
identify which decisions can change that expense. For this thesis, the focus
is the assignment and succession of product-carrying operations, and the
trade-off between their changeover costs and travel costs.

\section{Solution methods}\label{sec:solution-methods}

\subsection{Exact methods}

Exact methods provide a way to search for an optimal solution and, when the
search is completed with an appropriate bound, certify its quality.
Mathematical formulations are also useful for establishing reference
solutions and comparing alternative algorithms. The routing literature
contains both exact and heuristic treatments of multi-trip
problems~\cite{cattaruzza2016multipletrips}.

The selected application studies illustrate different exact approaches.
Lahyani et al.~\cite{lahyani2015oilmulticompartment} strengthen their integer
formulation with valid inequalities in a branch-and-cut algorithm.
Tamburini et al.~\cite{TAMBURINI2025104019} use column generation, with vessel
routes represented by columns. These approaches should be distinguished
from simply giving a formulation to a general-purpose solver: they exploit
the structure of the problem in different ways.

In this thesis, the MILP is used as a reference model. Its computational
results must be interpreted according to the solver's status and bounds.
A feasible solution returned at a time limit is a useful baseline, but is
not necessarily optimal. This distinction prevents a comparison against a
MILP incumbent from being presented as a comparison against a proven optimum.

\subsection{Heuristics and hybrid methods}

Laporte et al.~\cite{laporte2000heuristics} review classical and modern
heuristics for vehicle routing, including route-construction approaches and
tabu search. Their survey provides background for the distinction between
building a feasible plan and improving an existing plan. A heuristic can
make selected decisions quickly, but its output does not in general provide
an optimality certificate.

For rich problems, decomposition can separate decisions that are difficult
to optimize simultaneously. Miranda et al.~\cite{MIRANDA2018211} provide an
example in which production and routing subproblems exchange information
iteratively. The way this exchange is organized matters: fixing a decision
once and revisiting it during the search lead to different opportunities for
improvement. This is a useful distinction when describing the construction
stage developed for MPVRP-CC.

Sacramento, Solnon, and Pisinger~\cite{sacramento2020feeder} combine CP with
adaptive large neighbourhood search for feeder-vessel operations in
multi-terminal ports. CP is used to intensify the search when a new
best-known solution is found. Their study shows how a constraint model can
operate within a heuristic framework, rather than only as a standalone
solver. It concerns a port scheduling setting and does not supply a model
of product changeovers for the present problem.

These studies motivate examining how tasks are divided between a heuristic
and an optimization model. They do not imply that a particular combination
will always perform better. For MPVRP-CC, the construction choices and the
decisions left to CP must be made explicit, and the complete procedure must
be evaluated under a common computational budget.

\subsection{Constraint programming for routing}

CP has a long-standing role in routing research. Shaw~\cite{shaw1998constraint}
uses large neighbourhood search in which related visits are removed from a
plan and reinserted through constraint-based search. This work establishes
an early connection between routing improvement and CP. It also shows why
CP-based routing should not be presented as a new methodological category
introduced by this thesis.

Different CP representations emphasize different aspects of a route.
Vali and Salimifard~\cite{vali2017mtsp} model the multiple traveling salesman
problem using interval variables and global constraints. Their work offers
a scheduling-oriented representation of tours. Vismara and
Briot~\cite{vismara2018circuit} propose the \texttt{WeightedSubCircuits}
constraint, which represents several disjoint circuits covering selected
vertices under a total-cost bound. This is a circuit-based alternative,
with a different representation from the sequence variables used here.

Delecluse, Schaus, and Van Hentenryck~\cite{delecluse2022sequence} introduce
an insertion-based sequence-variable domain for routing. It supports search
on partial tours and constraints for capacities and time windows. Their
experiments include dial-a-ride, patient transportation, and the asymmetric
traveling salesman problem with time windows. Their later
preprint~\cite{delecluse2025sequence} formalizes the domain, its operations,
and consistency notions, and discusses optional visits and insertion-based
search.

For MPVRP-CC, the interest of this representation is that the order of
operations is itself a decision with an associated transition cost. However,
a suitable variable domain is only one part of a solution method. The model
must still express the relevant operational rules, account for costs
correctly, and provide an effective search strategy. Evidence that sequence
variables work well on another routing problem is motivation for this
investigation, not evidence of their performance on MPVRP-CC.

\section{Research gap}\label{sec:research-gap}

The reviewed literature establishes three foundations for this thesis.
First, repeated trips and multiple products are established routing
features. Second, cleaning and sequence-dependent preparation already occur
in transport and integrated production-routing studies. Third, CP has been
used both directly and within hybrid methods, with several representations
available for routing and sequencing.

Accordingly, the contribution sought here is not the introduction of
cleaning into routing in general, nor the first use of CP for routing. The
study focuses on a specified delivery problem in which product transitions
between successive loads contribute to the objective, together with travel
costs. It examines a solution approach based on CP and sequence variables
and evaluates the consequences of including those transition costs in
planning.

The selected papers provide relevant components and precedents, but the
comparisons above do not establish that their models and experiments answer
this precise combination of questions. This is a bounded research position:
it does not claim that no existing general formulation could express the
problem as a special case. The definition in
Section~\ref{sec:problem-definition} is therefore essential to understanding
what is being studied and to making subsequent comparisons meaningful.

Table~\ref{tab:research-gap} connects the literature to the questions pursued
in this thesis. Rather than interpreting unverified features as absent, it
states what the reviewed work contributes and what will be assessed here.

\begin{table}[htbp]
  \centering
  \small
  \caption{Relationship between the reviewed literature and this study}
  \label{tab:research-gap}
  \begin{tabularx}{\textwidth}{@{}XX@{}}
    \toprule
    Foundation in the literature & Question pursued in this thesis \\
    \midrule
    Multi-trip and rich routing models
      & How should the delivery, loading, and product-transition decisions
        be represented together? \\[0.5em]
    Cleaning charges and sequence-dependent preparation
      & How do routes, product sequences, and total costs change when
        changeovers enter the planning objective? \\[0.5em]
    CP routing representations and hybrid methods
      & How effectively does the proposed sequence-variable approach solve
        the specified problem compared with a MILP baseline? \\
    \bottomrule
  \end{tabularx}
\end{table}

The next chapter introduces the concepts needed to understand the models.
The Methodology chapter then defines MPVRP-CC and explains the approach
actually developed. In particular, any choices fixed before CP search must
be identified so that experimental conclusions refer to the decision space
that was explored.
